\clearpage\thispagestyle{empty}
\pdfbookmark[0]{Quelle zum Brief an von Lindenau vom 20. Mai 1809}{lindenau1809may-sources}
\begin{center}\textbf{Quelle zum Brief an von Lindenau vom 20. Mai 1809}\end{center}
\noindent Diese KI-Übertragung ins heutige Deutsch folgt dem Brief von Carl Friedrich Gauß an Bernhard August von Lindenau in \href{https://gdz.sub.uni-goettingen.de/id/PPN236007467}{Gauß’ \textit{Werken}, Bd.~VI} (Göttingen, 1874), S.~311--312. Der Brief erschien zuerst als Beitrag XLV, \textit{Auszug aus einem Schreiben des Herrn Prof. Gauß}, im Maiheft 1809 der \textit{Monatlichen Correspondenz}, Bd.~XIX, S.~\href{https://archive.org/download/bub_gb_caMAAAAAMAAJ/page/n515.jpg}{501}, \href{https://archive.org/download/bub_gb_caMAAAAAMAAJ/page/n516.jpg}{502} und \href{https://archive.org/download/bub_gb_caMAAAAAMAAJ/page/n517.jpg}{503}. Das digitalisierte Exemplar stammt von der New York Public Library.
\par\smallskip\noindent Die Werkausgabe fügt die beiden eckigen Klammern samt Inhalt hinzu. Sie betreffen von Lindenaus Veröffentlichung der Ceres-Ephemeride und die Anzeige von Gauß’ Programm zur Polhöhenbestimmung im Februarheft. Im Zitat zur Wahl des Halbkreises steht in beiden Drucken $\cos h$. KI-Übertragung und KI-Satz.
\clearpage
